% Arabic book typography. The two Amiri fonts are included in this project.
\usepackage{iftex}
\ifXeTeX\else
  \PackageError{BurtonKasidah}{Please compile this project with XeLaTeX}
    {In Overleaf, choose XeLaTeX as the compiler and main.tex as the main document.}
\fi
\usepackage{fontspec}
\usepackage[a4paper,top=23mm,bottom=24mm,left=23mm,right=23mm]{geometry}
\usepackage{setspace}
\usepackage{needspace}
\usepackage{parskip}
\usepackage[unicode,pdfencoding=auto]{hyperref}
\newenvironment{SourceQuotation}{%
  \begin{quote}
  \setlength{\parskip}{0pt}%
}{\end{quote}}
% Allow a recent bidi package to run with an older LaTeX kernel.
\makeatletter
\providecommand\IfClassLoadedT[2]{\@ifclassloaded{#1}{#2}{}}
\providecommand\IfPackageLoadedT[2]{\@ifpackageloaded{#1}{#2}{}}
\providecommand\IfPackageLoadedF[2]{\@ifpackageloaded{#1}{}{#2}}
\providecommand\IfFileLoadedTF[1]{%
  \expandafter\ifx\csname ver@#1\endcsname\relax
    \expandafter\@secondoftwo
  \else
    \expandafter\@firstoftwo
  \fi}
\makeatother
% Load polyglossia after the packages whose RTL behaviour bidi adapts.
\usepackage{polyglossia}
\setdefaultlanguage[numerals=mashriq]{arabic}
\setotherlanguage{english}

\setmainfont[
  Path=fonts/,
  BoldFont=Amiri-Bold.ttf,
  ItalicFont=Amiri-Regular.ttf,
  BoldItalicFont=Amiri-Bold.ttf,
  ItalicFeatures={FakeSlant=0.15},
  BoldItalicFeatures={FakeSlant=0.15}
]{Amiri-Regular.ttf}
\newfontfamily\arabicfont[
  Path=fonts/,
  Script=Arabic,
  Language=Arabic,
  BoldFont=Amiri-Bold.ttf,
  ItalicFont=Amiri-Regular.ttf,
  BoldItalicFont=Amiri-Bold.ttf,
  ItalicFeatures={FakeSlant=0.15},
  BoldItalicFeatures={FakeSlant=0.15}
]{Amiri-Regular.ttf}
\newfontfamily\englishfont[
  Path=fonts/,
  BoldFont=Amiri-Bold.ttf,
  ItalicFont=Amiri-Regular.ttf,
  BoldItalicFont=Amiri-Bold.ttf,
  ItalicFeatures={FakeSlant=0.15},
  BoldItalicFeatures={FakeSlant=0.15}
]{Amiri-Regular.ttf}

\hypersetup{
  pdftitle={قصيدة الحاج عبدو اليزدي: ترجمة عربية أولية},
  pdfauthor={Richard Francis Burton; preliminary Arabic translation by ChatGPT / OpenAI},
  pdfsubject={Arabic preliminary translation with historical and editorial notes},
  pdfdisplaydoctitle=true,
  colorlinks=true,
  linkcolor=black,
  urlcolor=black,
  bookmarksopen=true,
  bookmarksnumbered=false
}
\setstretch{1.22}
\setlength{\parindent}{0pt}
\setlength{\parskip}{0.65\baselineskip plus 2pt minus 1pt}
\setlength{\emergencystretch}{2em}
\tolerance=1600
\widowpenalty=10000
\clubpenalty=10000
\raggedbottom
\setcounter{tocdepth}{1}
\renewcommand{\contentsname}{محتويات الكتاب}

% Each substantial source section begins on a fresh page.
\newcommand{\MajorSection}[1]{%
  \clearpage
  \phantomsection
  \section*{#1}
  \addcontentsline{toc}{section}{#1}%
}
\newcommand{\MinorSection}[1]{%
  \Needspace{5\baselineskip}
  \subsection*{#1}%
}

% A complete couplet and any following Burton gloss stay together.
% Arguments: Eastern Arabic number, first line, second line, optional gloss.
\newcommand{\Couplet}[4]{%
  \par\noindent
  \begin{minipage}{\linewidth}
    \setlength{\parskip}{0pt}%
    \setstretch{1.25}%
    {\small\bfseries #1}\enspace #2\par
    #3\par
    #4%
  \end{minipage}%
  \par\vspace{0.45\baselineskip}%
}
\newcommand{\BurtonNote}[1]{%
  \vspace{0.3\baselineskip}%
  {\small\setstretch{1.15} #1\par}%
}

% Values recorded by the project initiator for the FIRST Arabic draft.
% They do not describe the later LaTeX conversion session.
\newcommand{\RecordedModel}{\textenglish{GPT-6.1 Sol Max}}
\newcommand{\RecordedDraftTime}{\textenglish{1h 10m 13s}}
